\chapter{Logic}
\section{Logical Connectives and Quantifiers}
\begin{definition}
    If $p$ and $q$ are propositions, then
    \begin{enumerate}[itemindent=0ex, itemsep=0pt, label=\upshape{(\alph*)}]
        \item the \emph{conjunction} of $p$ and $q$, denoted by $p \wedge q$, is the proposition ``$p$ and $q$'', which is true if both $p$ and $q$ are true, and false otherwise.
        \item the \emph{disjunction} of $p$ and $q$, denoted by $p \vee q$, is the proposition ``$p$ or $q$'', which is true if at least one of $p$ and $q$ is true, and false otherwise.
        \item the \emph{negation} of $p$, denoted by $\neg p$, is the proposition ``not $p$'', which is true if $p$ is false, and false if $p$ is true.
    \end{enumerate}
\end{definition}

\begin{example}
    Set $p=$ ``It is raining'' and $q=$ ``I have an umbrella''. Then $p \wedge q$ is the proposition ``It is raining and I have an umbrella'', $p \vee q$ is the proposition ``It is raining or I have an umbrella'', and $\neg p$ is the proposition ``It is not raining''.
\end{example}

\begin{definition}
    We define two quantifiers as follows:
    \begin{enumerate}[itemindent=0ex, itemsep=0pt, label=\upshape{(\alph*)}]
        \item the \emph{universal quantifier} $\forall$ is used to denote ``for all'' or ``for every''.
        \item the \emph{existential quantifier} $\exists$ is used to denote ``there exists'' or ``for some''.
    \end{enumerate}
\end{definition}

We can combine the quantifiers with logical connectives to form more complex propositions. 

\begin{example}
    Set $E(t)=$ ``$t$ is even'' and $P(t)=$ ``$t$ is prime''. Then $\exists t\quad (E(t) \wedge P(t))$ is the proposition ``There exists a $t$ such that $t$ is even and $t$ is prime''.
\end{example}

\begin{definition}
    The \emph{conditional statement} $p \rightarrow q$ is the proposition ``if $p$ then $q$'', which is false if $p$ is true and $q$ is false, and true otherwise.
\end{definition}

$p \rightarrow q$ can also be read as ``$p$ implies $q$'', hence also called an \emph{implication}. Usually it is used after the universal quantifier $\forall$. For example, $\forall t\quad (S(t) \rightarrow C(t))$ is the proposition ``For all $t$, if $S(t)$ then $C(t)$''. 

Now that $\forall$ is associated with $\rightarrow$, what about $\exists$? In fact, $\exists$ is associated with $\wedge$. For example, $\exists t\quad (U(t) \wedge M(t))$ is the proposition ``There exists a $t$ such that $U(t)$ and $M(t)$''. If we set $U(t)=$ ``$t$ is a US student'' and $M(t)=$ ``$t$ has visited Mexico'', then $\exists t\quad (U(t) \wedge M(t))$ is the proposition ``There exists a US student who has visited Mexico''.

We now consider some rather complex examples. 

\begin{example}
    The sentence in natural language ``For every student in SAI, a research mentor will be assigned to him or her'' can be translated into the following proposition:
    \[\forall x\quad (S(x) \rightarrow \exists y\quad M(x, y))\]
    where $S(t)=$ ``$t$ is a student in SAI'' and $M(t_1, t_2)=$ ``$t_2$ is a research mentor assigned to $t_1$''.

    The sentence in natural language ``Even if there is only one student in SAI who does not study hard, he or she will get a very bad score'' can be translated into the following proposition:
    \[\forall x\quad (S(x) \wedge \neg H(x)) \rightarrow B(x)\]
    where $S(t)=$ ``$t$ is a student in SAI'', $H(t)=$ ``$t$ studies hard'', and $B(t)=$ ``$t$ gets a very bad score''. Note that the sentence in natural language contains the phrase ``Even if'', which is a bit tricky. We can rewrite the sentence as ``For every student in SAI, if he or she does not study hard, then he or she will get a very bad score''. That's why we used the universal quantifier $\forall$ and the implication $\rightarrow$.
\end{example}

\begin{example}
    Let's consider how to express the classic mathematical statement ``There is a unique $x$ such that $P(x)$'' using quantifiers and logical connectives. We can express it as follows:
    \[\exists x\quad (P(x) \wedge \forall y\quad (P(y) \rightarrow y = x))\]
    This proposition states that there exists an $x$ such that $P(x)$ is true, and for all $y$, if $P(y)$ is true, then $y$ must be equal to $x$. In other words, there is exactly one $x$ that satisfies the property $P$.
\end{example}

\section{Inference}
We start with a simple example. Suppose we have the following propositions:
\begin{align*}
    p\colon &T\text{ is connected and has no simple circuits.} \\
    q\colon &T\text{ is connected.}
\end{align*}
From $p$, we can infer $q$. In other words, if $p$ is true, then $q$ must be true. We can write this inference in symbols: $p\wedge q \vdash q$. Note that in this example, we don't have to know the exact definition of $T$ being connected or having no simple circuits. We only need to know that $p$ implies $q$. 

There is a special example: $\vdash p \vee \neg p$. No assumption is needed in this example. 

We now define the new symbol $\vdash$ formally.

\begin{definition}
    If $\Phi$ is a set of propositions and $\psi$ is a proposition, then we write $\Phi \vdash \psi$ to represent that \emph{$\psi$ is derivable from $\Phi$}. $\vdash\psi$ is the shorthand for $\emptyset \vdash \psi$, and $\Phi, \phi_1, \dots, \phi_n \vdash \psi$ is the shorthand for $\Phi \cup \{\phi_1, \dots, \phi_n\} \vdash \psi$.
\end{definition}

We can extract some common patterns from concrete examples. Here is a summary. 

\begin{example}
    We have the following inference rules, the so-called ``single-step proofs'':
    \begin{enumerate}[itemindent=0ex, itemsep=0pt, label=\upshape{(\alph*)}]
        \item $p\wedge q\vdash p$. 
        \item $p\vdash p\vee q$.
        \item $\vdash p\vee \neg p$.
        \item $\forall x.\ M(x)\vdash M(c)$.
        \item $M(c)\vdash \exists x.\ M(x)$.
    \end{enumerate}
\end{example}

Now let's take a look at some more complicated examples. 

\begin{example}
    Consider the proof of the following proposition: 
    \[\forall\varepsilon>0,\ \exists N, \forall n>N,\ \left|\frac{1}{n}\right|<\varepsilon.\]

    \begin{proof}
        Given $\varepsilon>0$, let $N=\left[ \frac{1}{\varepsilon} \right]+1$. Then for all $n>N$, we have $\left|\frac{1}{n}\right|=\frac{1}{n}<\frac{1}{N}<\varepsilon$. 
    \end{proof}

    Proof step: In order to show $\forall\varepsilon>0,\ \exists N, \forall n>N,\ \left|\frac{1}{n}\right|<\varepsilon$, it suffices to show that if $\varepsilon>0$ then $\exists N, \forall n>N,\ \left|\frac{1}{n}\right|<\varepsilon$.

    The logic behind this step: If $\Phi, P(x)\vdash Q(x)$ then $\Phi\vdash \forall x.\ (P(x)\rightarrow Q(x))$. 
    
    \mbox{}

    Proof step: In order to show $\exists N, \forall n>N,\ \left|\frac{1}{n}\right|<\varepsilon$, it suffices to prove $\forall n>\left[ \frac{1}{\varepsilon} \right]+1,\ \left|\frac{1}{n}\right|<\varepsilon$.

    The logic behind this step: If $\Phi\vdash P(t)$ then $\Phi\vdash \exists x.\ P(x)$.
\end{example}

\begin{example}
    The following logics are used in the proof of irrationality of $\sqrt{2}$:
    \begin{enumerate}[itemindent=0ex, itemsep=0pt, label=\upshape{(\alph*)}]
        \item If $\Phi$ does not talk about $x$ and $y$, $\psi_1$ does not talk about $y$, and $\Phi, \psi_1, \psi_2, \psi_3\vdash\neg\chi$, then $\Phi\vdash \neg\exists x(\psi_1\wedge\exists y(\psi_2\wedge\psi_3\wedge\chi))$.
        \item If $\Phi\vdash\psi$ and $\Phi,\psi\vdash\chi$, then $\Phi\vdash\chi$.
        \item If $\Phi,P(x)\vdash\chi$, then $\Phi,\exists x.\ P(x)\vdash\chi$.
    \end{enumerate}
\end{example}

We now dive into propositional logic. 

\begin{definition}
    We define the symbols \textsc{assumption}, \textsc{weakening}, and \textsc{substitution} as follows:
    \begin{enumerate}[itemindent=0ex, itemsep=0pt, label=\upshape{(\alph*)}]
        \item \textsc{assumption}. If $\varphi\in\Phi$, then $\Phi\vdash\varphi$.
        \item \textsc{weakening}. If $\Phi\subseteq\Phi'$ and $\Phi\vdash\varphi$, then $\Phi'\vdash\varphi$.
        \item \textsc{substitution}. If $\Phi\vdash\psi$ and $\Phi,\psi\vdash\chi$, then $\Phi\vdash\chi$.
    \end{enumerate}
\end{definition}

\begin{example}
    Let's prove the following derivations using \textsc{assumption}, \textsc{weakening}, and \textsc{substitution}. 
    \begin{enumerate}[itemindent=0ex, itemsep=0pt, label=\upshape{(\alph*)}]
        \item $\varphi\vdash\varphi$.
        \item If $\varphi\vdash\psi$ and $\psi\vdash\chi$, then $\varphi\vdash\chi$.
        \item If $\Phi, \varphi\vdash\psi$ and $\Phi, \psi\vdash\chi$, then $\Phi, \varphi\vdash\chi$.
    \end{enumerate}
\end{example}

\begin{proof}
    \begin{enumerate}[itemindent=0ex, itemsep=0pt, label=\upshape{(\alph*)}]
        \item $\varphi\vdash\varphi$ (by \textsc{assumption}).
        \item (1) $\varphi\vdash\psi$ (by hypothesis). 
        
        (2) $\psi\vdash\chi$ (by hypothesis). 
        
        (3) $\varphi, \psi\vdash\chi$ (by \textsc{weakening} and (2)). 
        
        (4) $\varphi\vdash\chi$ (by \textsc{substitution} and (1) and (3)).
        \item (1) $\Phi, \varphi\vdash\psi$ (by hypothesis). 
        
        (2) $\Phi, \psi\vdash\chi$ (by hypothesis). 
        
        (3) $\Phi, \varphi, \psi\vdash\chi$ (by \textsc{weakening} and (2)). 
        
        (4) $\Phi, \varphi\vdash\chi$ (by \textsc{substitution} and (1) and (3)).
    \end{enumerate}
\end{proof}

\begin{definition}
    We define the symbols $\textsc{and-intro}$, $\textsc{and-elim1}$, and $\textsc{and-elim2}$ as follows:
    \begin{enumerate}[itemindent=0ex, itemsep=0pt, label=\upshape{(\alph*)}]
        \item \textsc{and-intro}. $\varphi,\psi\vdash\varphi\wedge\psi$. 
        \item \textsc{and-elim1}. $\varphi\wedge\psi\vdash\varphi$.
        \item \textsc{and-elim2}. $\varphi\wedge\psi\vdash\psi$.
    \end{enumerate}
\end{definition}

\begin{example}
    Let's prove the following derivations. 
    \begin{enumerate}[itemindent=0ex, itemsep=0pt, label=\upshape{(\alph*)}]
        \item If $\Phi\vdash\psi$ and $\Phi\vdash\chi$, then $\Phi\vdash\psi\wedge\chi$.
        \item If $\Phi\vdash\psi\wedge\chi$, then $\Phi\vdash\psi$ and $\Phi\vdash\chi$.
        \item If $\Phi,\psi_1\wedge\psi_2\vdash\chi$, then $\Phi,\psi_1,\psi_2\vdash\chi$.
        \item If $\Phi, \psi_1\vdash\chi$, then $\Phi, \psi_1\wedge\psi_2\vdash\chi$.
        \item If $\Phi, \psi_2\vdash\chi$, then $\Phi, \psi_1\wedge\psi_2\vdash\chi$.
    \end{enumerate}
\end{example}

\begin{proof}
    \begin{enumerate}[itemindent=0ex, itemsep=0pt, label=\upshape{(\alph*)}]
        \item (1) $\psi,\chi\vdash\psi\wedge\chi$ (by \textsc{and-intro}). 
        
        (2) $\Phi,\psi,\chi\vdash\psi\wedge\chi$ (by \textsc{weakening} and (1)). 

        (3) $\Phi\vdash\chi$ (by hypothesis). 
        
        (4) $\Phi,\psi\vdash\chi$ (by \textsc{weakening} and (3)). 
        
        (5) $\Phi,\psi\vdash\psi\wedge\chi$ (by \textsc{substitution} and (4) and (2)). 
        
        (6) $\Phi\vdash\psi$ (by hypothesis). 
        
        (7) $\Phi\vdash\psi\wedge\chi$ (by \textsc{substitution} and (6) and (5)).
        \item (1) $\Phi\vdash\psi\wedge\chi$ (by hypothesis). 
        
        (2) $\psi\wedge\chi\vdash\psi$ (by \textsc{and-elim1}). 
        
        (3) $\Phi, \psi\wedge\chi\vdash\psi$ (by \textsc{weakening} and (2)). 
        
        (4) $\Phi\vdash\psi$ (by \textsc{substitution} and (1) and (3)). 
        
        The proof of $\Phi\vdash\chi$ is similar.
        \item (1) $\psi_1,\psi_2\vdash\psi_1\wedge\psi_2$ (by \textsc{and-intro}). 
        
        (2) $\Phi,\psi_1,\psi_2\vdash\psi_1\wedge\psi_2$ (by \textsc{weakening} and (1)). 
        
        (3) $\Phi, \psi_1\wedge\psi_2\vdash\chi$ (by hypothesis). 
        
        (4) $\Phi,\psi_1,\psi_2,\psi_1\wedge\psi_2\vdash\chi$ (by \textsc{weakening} and (3)). 
        
        (5) $\Phi,\psi_1,\psi_2\vdash\chi$ (by \textsc{substitution} and (2) and (4)).
    \end{enumerate}
\end{proof}